\chapter*{Tóm tắt}
\addcontentsline{toc}{chapter}{Tóm tắt}

Trầm cảm là một trong những rối loạn sức khỏe tâm thần phổ biến nhất trên toàn cầu, gây suy giảm nghiêm trọng chất lượng cuộc sống và tạo gánh nặng kinh tế - xã hội to lớn. Việc đánh giá mức độ nghiêm trọng của triệu chứng trầm cảm hiện nay chủ yếu dựa vào các bảng câu hỏi tự điền, điển hình là thang đo Patient Health Questionnaire-9 (PHQ-9) với tổng điểm từ 0 đến 27. Mặc dù các mô hình học máy (Machine Learning) nhị phân ($T_2$: Không có triệu chứng vs Có triệu chứng) đã đạt được những kết quả khả quan, hiệu năng dự đoán bị suy giảm nghiêm trọng khi tăng độ chi tiết của nhãn sang phân loại đa mức ($T_3, T_4, T_5$), đặc biệt là sự khó khăn trong việc phân biệt các mức độ triệu chứng kế cận.

Công trình này nghiên cứu, xây dựng và đánh giá toàn diện khung học máy \textbf{DepSev-ML} (Depression Severity Machine Learning) nhằm dự đoán mức độ triệu chứng trầm cảm dựa trên dữ liệu sinh lý, cận lâm sàng, lối sống và nhân khẩu học từ Khảo sát Kiểm tra Sức khỏe và Dinh dưỡng Quốc gia Hoa Kỳ (NHANES). Nghiên cứu được triển khai trên hai nhánh dữ liệu chính: Legacy Development Cohort (chu kỳ 2013--2014, $N = 4.045$, $F_4 = 34$ đặc trưng) và Official Multi-cycle Cohort (các chu kỳ 2011--2016, $N = 14.581$, $F_{4\_\text{common31}} = 31$ đặc trưng hài hòa).

Nghiên cứu khảo sát ba trường phái phân loại chính: Phân loại phẳng (Flat Classification $E_0$), Phân loại phân cấp hai giai đoạn (Hierarchical Classification gồm Hard Routing $E_1$--$E_3b$, Soft Probability Fusion $E_4$--$E_5$ và Cost-Sensitive Decision Making $E_6a$--$E_6b$), và Phân loại thứ bậc (Ordinal Classification $E_7a$--$E_7b$). 

Kết quả thực nghiệm chỉ ra những phát hiện khoa học then chốt: (1) Trên không gian $T_5$ (5 mức), Macro-F1 suy giảm từ 0.6639 ($T_2$) xuống 0.2630 ($T_5$), trong đó bước sụt giảm lớn nhất xảy ra khi triệu chứng nhẹ tách thành lớp độc lập; việc mở rộng dữ liệu đa chu kỳ ($N = 14.581$) không giúp cải thiện Macro-F1 nhất quán mà chỉ tăng nhẹ Balanced Accuracy từ 0.2841 lên 0.3076. (2) Trên không gian $T_3$ (3 mức), mô hình Hard Hierarchical ($E_1$) cải thiện Macro-F1 lên 0.4550 nhưng gây tắc nghẽn nghiêm trọng khi loại bỏ nhầm 120 ca nặng quy ước ($C_2^+$) ở Giai đoạn 1. (3) Mô hình Soft Hierarchical kết hợp tối ưu chi phí Bayes ($E_6b$) giảm số ca bỏ sót $C_2^+ \to C_0$ xuống chỉ còn 62 ca. (4) Mô hình Ordinal Logistic Regression cân bằng ($E_7a$) thiết lập đỉnh cao mới về Balanced Accuracy (0.5151) và đạt Recall $C_2^+ = 60.37\%$, duy trì độ bền vững trên tập kiểm tra độc lập theo dòng thời gian ($N = 4.888$, Balanced Accuracy = 0.5073). Tuy nhiên, nghiên cứu cũng chỉ ra rằng việc tối ưu hóa sai số tuyệt đối thứ bậc (Class MAE) trên dữ liệu mất cân bằng nặng sẽ gây suy biến lớp trung gian ($C_1$), khẳng định sự cần thiết của việc đánh giá đa chỉ số khách quan.

\vspace{0.8cm}
\noindent \textbf{Từ khóa:} Trầm cảm, PHQ-9, Dữ liệu sinh lý, Học máy, Phân loại phân cấp, Phân loại thứ bậc, Quyết định nhạy chi phí, NHANES, DepSev-ML.
\cleardoublepage
